\section{Validación y resultados disponibles}
\subsection{Demostración local del 27 de octubre}
Se ejecutaron 32 pruebas Python y 8 Node, todas aprobadas. La integración
PostgreSQL 16.6/MongoDB 7.0.16 terminó en PASS con tres ejemplos publicados,
uno de aire, uno meteorológico y uno de tráfico, adaptados en copias de prueba.
Los ejemplos de meteo/tráfico en staging no son registros aceptados de la carga anual.

La integración comprobó restricciones SQL/BSON, índices, consultas, replay,
conflictos en cuarentena, CRUD aislado y respaldo/restauración lógico de ambos
motores. También probó ocho conexiones PostgreSQL sobre una clave y dos payloads:
una aceptación, tres duplicados y cuatro cuarentenas. No es estrés sostenido ni
prueba de concurrencia MongoDB entre procesos \cite{integracion,respaldos}.

Los contenedores fueron locales, sin puertos publicados y con red deshabilitada.
Se eliminaron al finalizar. El informe generado fue
\archivo{/tmp/opencode/smartcity-s1-integration-t7d2ibo2.json}; es un artefacto
local, no un archivo disponible en cualquier clon.

\subsection{Carga anual: evidencia previa, no reejecución de hoy}
El mínimo es de 100.000 registros PostgreSQL y 500.000 documentos MongoDB
según el enunciado y el plan del equipo \cite{enunciado,plan}.
\begin{longtable}{p{0.48\textwidth}p{0.2\textwidth}p{0.2\textwidth}}
\toprule Control registrado & PostgreSQL & MongoDB \\
\midrule\endhead
Lecturas/documentos aceptados & 1.076.813 & 1.076.813 \\
Metadatos y valores comparados & 1.076.813 & 1.076.813 \\
Replay del lote completo & Sin aumento & Sin aumento \\
Mediciones sintéticas & 0 & 0 \\
Mínimo de volumen del motor & Superado localmente & Superado localmente \\
Segunda fuente de mediciones aceptada & No & No \\
\bottomrule
\end{longtable}

Estos resultados proceden de los manifiestos
\archivo{research/datasets/results/expanded-2025/postgresql_full.json} y
\archivo{mongodb_full.json}, y del cierre documentado \cite{cierre}.
La revisión actual contrastó los hashes de sus cuatro scripts registrados con
el código disponible; no repitió la carga masiva ni verificó todos los insumos privados.

Los mínimos locales no prueban que el servidor compartido tenga esos datos.
Tampoco constituyen una comparación de rendimiento entre motores: almacenamiento,
configuración y operaciones deben ser equivalentes antes de comparar tiempos.

\subsection{Reproducción del núcleo}
Desde la raíz del repositorio, con Python 3.10+ y Node.js 18+:
\begin{lstlisting}
bash scripts/validate_sprint1.sh
\end{lstlisting}
Para integrar motores se requiere Docker local y las imágenes ya descargadas
\archivo{postgres:16.6-alpine} y \archivo{mongo:7.0.16}:
\begin{lstlisting}
PYTHONDONTWRITEBYTECODE=1 python3 scripts/test_sprint1_engines.py
\end{lstlisting}
El harness no descarga imágenes. Las instrucciones completas y su alcance están
en \archivo{ENTREGA_SPRINT1.md}. Los clientes básicos psycopg/PyMongo aún
necesitan prueba de extremo a extremo con autenticación real.

\subsection{Respaldo y restauración}
La prueba usa pg\_dump/pg\_restore y mongodump/mongorestore, restaura en otra
base desechable y compara contenido y estructura. Esto acredita el caso
controlado, no una política de respaldo productiva, recuperación tras caída,
retención externa o restauración del servidor compartido.

\subsection{Qué evidencia falta adjuntar}
\begin{itemize}
\item Salidas verificables de cada una de las veinte consultas en la ruta elegida.
\item Diagrama E--R completo y dependencias funcionales revisados por Daniel.
\item Modelo documental final, ejemplos BSON y revisión de Kembly.
\item Capturas legibles y registro de revisión cruzada, sin credenciales ni hosts privados.
\item Comprobación de acceso, conteos y respaldo compartidos si se usan en la entrega.
\end{itemize}
